\documentclass[11pt,letterpaper]{article}
\usepackage[T1]{fontenc}
\usepackage{lmodern}
\usepackage[margin=1in,headheight=15pt]{geometry}
\usepackage{amsmath,amssymb,amsthm,mathtools}
\usepackage{microtype,booktabs,longtable,array,xcolor,xurl}
\usepackage{fancyhdr}
\usepackage[colorlinks=true,linkcolor=blue!40!black,citecolor=blue!40!black,urlcolor=blue!50!black]{hyperref}
\usepackage{bookmark}
\allowdisplaybreaks[2]
\emergencystretch=2em
\setlength{\parskip}{3pt}
\numberwithin{equation}{section}
\newtheorem{theorem}{Theorem}[section]
\newtheorem{proposition}[theorem]{Proposition}
\newtheorem{lemma}[theorem]{Lemma}
\newtheorem{corollary}[theorem]{Corollary}
\theoremstyle{definition}
\newtheorem{definition}[theorem]{Definition}
\theoremstyle{remark}
\newtheorem{remark}[theorem]{Remark}
\DeclareMathOperator{\Li}{Li}
\DeclareMathOperator{\Rea}{Re}
\DeclareMathOperator{\Res}{Res}
\newcommand{\Rcal}{\mathcal R}
\newcommand{\Pcal}{\mathcal P}
\newcommand{\Ical}{\mathfrak I}
\newcommand{\dd}{\,\mathrm d}
\newcommand{\file}[1]{\nolinkurl{#1}}
\newcommand{\pin}{\texttt{16c7e342d7a4a15f59911f322eb90b7f96c2a8b5}}
\pagestyle{fancy}
\fancyhf{}
\fancyhead[L]{\small Centered Dougall resonance}
\fancyhead[R]{\small ProveIt research continuation}
\fancyfoot[C]{\thepage}
\title{\textbf{Centered Dougall Resonance}\\[7pt]
\large Exact Stieltjes Jets, Harmonic Identities,\\
Polylogarithmic Subtraction, and Zeta Antiderivatives}
\author{Prepared for Vladimir Reshetnikov's ProveIt programme\\
\small Mathematical derivations and computational audit with OpenAI assistance}
\date{10 October 2026}
\hypersetup{pdftitle={Centered Dougall Resonance},pdfauthor={ProveIt research continuation with OpenAI assistance},pdfsubject={Exact identities for Rogers-Dougall sums, Hurwitz zeta and Stieltjes constants}}
\begin{document}
\maketitle
\thispagestyle{plain}
\begin{abstract}
The incoming Gauss--Hurwitz article asks for a Dougall-type summation with an
explicit residue polynomial and complete resonant spectral jets. We supply
such a construction for the Rogers--Dougall ${}_5F_4(1)$ sum. Centering the
summation variable at half the very-well-poised parameter removes every odd
inverse-power correction. The spectral exponent is doubled, so integral
and half-integral parameter values lead to two different zeta towers.

At every nonpositive integer parameter, an ordinary, explicitly subtracted
series collapses to one short digamma expression. Twice the Gamma-product
residue is
\[
 \frac{(-1)^N}{N!}(d_0)_N(1-b)_N(1-c)_N.
\]
The decisive ingredient is a differential identity for the Bernoulli
coefficient polynomials. It cancels all lower Hurwitz counterterms, rather
than merely leaving a finite but lengthy expression. Every higher spectral
jet has an explicit finite formula in Gamma/polygamma data, generalized
Stieltjes constants, and Hurwitz-zeta derivatives at negative odd integers.
The interlaced half-integer values instead involve nonpositive even
integers and no Stieltjes pole.

We derive all-order harmonic identities, including a first-resonance
identity containing $\zeta'(-1)$, and a continuous-shift harmonic formula.
Every Euler primitive of the first half-parameter harmonic jet has a closed
zeta endpoint. A separate finite calculus gives antiderivatives of all
collapsed resonance values in the centering parameter. The proofs are
accompanied by 485 exact assertions and 293 numerical diagnostics at 70
working decimal digits. Classical summation, parameter-differentiation,
and negative-polygamma antecedents are credited explicitly; no global
priority or proof-assistant formalization is claimed.
\end{abstract}
\clearpage
\begingroup
\small
\setlength{\parskip}{0pt}
\tableofcontents
\endgroup
\clearpage

\section{Research target, provenance, and contribution}
\label{dr:sec:scope}

The canonical ProveIt manuscript connects polylogarithms, harmonic sums,
Gamma functions, and Hurwitz-zeta jets. Its incoming article
\emph{Gauss--Hurwitz Pole Cancellation} makes a concrete proposal in its
Section~9, question~1: pass from Gauss summation to a fixed Dougall-type
family whose resonance residues and complete jets can be made explicit
\cite{incoming}. This article answers that specific proposal. It does not
classify all generalized hypergeometric summations admitting such a
calculus.

The principal result is not a new version of Dougall's classical summation
formula. It is the explicit centered subtraction and cancellation theorem
below, together with its derivative and antiderivative consequences. There
are three structural differences from the Gauss construction. The natural
large-$n$ coordinate is $n+\kappa$; its inverse-power expansion has only even
steps; and the exponent is $-1-2u$, not $-1-u$. Consequently the Stieltjes
pole has residue $1/2$ as a function of $u$, while negative half-integer
values of $u$ are regular points of the Gamma product.

The central proof is short once its coefficient differential law is found.
The surrounding analytic work is nevertheless necessary: the subtracted
sum must converge normally, all parameter derivatives must be legitimate,
and reciprocal-Gamma zeros must be preserved before differentiation. The
article treats each of these points directly. Its conclusions are ordinary
analytic identities, not deductions from small numerical residuals.

\subsection{Classical inputs and the priority boundary}

The Rogers--Dougall very-well-poised ${}_5F_4(1)$ identity is classical and is
recorded as DLMF~16.4.9 \cite{dlmf-dougall}. Bernoulli expansions for Gamma
ratios, Hurwitz-zeta continuation, and the Gamma--zeta derivative relation
are also classical \cite{dlmf-gamma,dlmf-hurwitz}. B\"uhring's work treats the
asymptotic behavior of generalized hypergeometric partial sums, including
the constant-term setting relevant to resonance \cite{buhring}. Thus a
finite-part viewpoint on divergent hypergeometric sums is not being
introduced here.

Parameter differentiation of hypergeometric identities is an established
source of harmonic-number identities. In particular, Chu and Fu treat
Dougall--Dixon formulas and harmonic identities \cite{chufu}; Nicholson
gives another account of hypergeometric transformations with harmonic
numbers \cite{nicholson}. The polynomial-weighted primitive construction
used later belongs to the Hurwitz-zeta and negative-polygamma framework of
Espinosa and Moll \cite{espinosa}, already present in the canonical
integration chapter \cite{repo-integration}.

The contribution supplied for integration is therefore specific: a
parity-adapted coefficient system, its differential cancellation law, the
all-integer collapsed formula, a degeneracy-safe all-jet identity, the
interlaced even-argument tower, and their explicitly proved specializations.
The limited literature inspection does not establish that each individual
harmonic specialization is absent from earlier work.

\subsection{Source scope and unresolved questions}

The manuscript README and integration chapter were inspected, together
with the incoming-directory listing and the correspondingly named Gauss--Hurwitz PDF
in the Library. The latter's question on printed page~19 is the target
addressed here. The last observed repository head was \pin. The source
reads were observations of a changing branch, not an exhaustive audit of
that entire commit. Individual file and archive identifiers are recorded
in \file{notes/SOURCE_AUDIT.md}.

The canonical README states that $S_4$ has an exact proof while $S_6$ and
the revised $S_8$ candidate remain conjectural \cite{repo-readme}. Nothing
in this article supplies their missing colored-polylogarithm certificates.
No numerical independence claim, new repository commit, or pull request is
made. The new results answer one expressly proposed research direction;
they do not change unrelated conjecture statuses.

\section{The normalized Rogers--Dougall family}
\label{dr:sec:family}

\subsection{Parameters, powers, and spectral derivatives}

Throughout the main construction, fix real parameters
\begin{equation}
 \kappa>0,\qquad b>0,\qquad c>0,\qquad
 d_0=1+2\kappa-b-c>0,
 \qquad x_n=n+\kappa.
 \label{dr:eq:domain}
\end{equation}
The variable $u$ is complex. Powers of $x_n$ use its real logarithm.
A prime or a superscript on $\zeta(s,a)$ differentiates its first,
\emph{spectral}, variable unless a parameter derivative is displayed.
We use
\begin{equation}
 \zeta(1+v,a)=\frac1v+
 \sum_{m=0}^{\infty}\frac{(-1)^m}{m!}\gamma_m(a)v^m,
 \qquad \gamma_0(a)=-\psi(a).
 \label{dr:eq:stieltjes}
\end{equation}
The ordinary Stieltjes constants are $\gamma_m=\gamma_m(1)$, with
$\gamma_0=\gamma$ the Euler constant. The notation $(a)_N$ denotes the
rising factorial, including $(a)_0=1$.

Define the summand by
\begin{align}
 W_n(u)={}&x_n
 \frac{\Gamma(n+2\kappa)\Gamma(n+b)}
 {\Gamma(n+1)\Gamma(n+1+2\kappa-b)}
 \frac{\Gamma(n+c)\Gamma(n+d_0-u)}
 {\Gamma(n+1+2\kappa-c)\Gamma(n+b+c+u)}.
 \label{dr:eq:weight}
\end{align}
Every reciprocal Gamma factor is interpreted as the entire function
$1/\Gamma$. In particular, a summand is allowed to be zero. On
$\Rea u<d_0$ all numerator Gamma arguments in \eqref{dr:eq:weight}
have positive real part, so there are no numerator poles. The fixed
denominator arguments are positive by \eqref{dr:eq:domain}.

\subsection{The classical summation in the required normalization}

Put
\begin{equation}
 P(u)=\frac{\Gamma(b)\Gamma(c)\Gamma(d_0-u)}
 {\Gamma(d_0)\Gamma(b+u)\Gamma(c+u)},
 \qquad G(u)=\frac12\Gamma(u)P(u).
 \label{dr:eq:gamma-side}
\end{equation}
Then the Rogers--Dougall formula gives
\begin{equation}
 \sum_{n=0}^{\infty}W_n(u)=G(u),
 \qquad 0<\Rea u<d_0.
 \label{dr:eq:dougall}
\end{equation}
For clarity, its standard form is
\begin{align}
 &{}_5F_4\!\left(\begin{matrix}a,1+a/2,b,c,d\\
 a/2,1+a-b,1+a-c,1+a-d\end{matrix};1\right)\notag\\
 &\qquad=
 \frac{\Gamma(1+a-b)\Gamma(1+a-c)\Gamma(1+a-d)
       \Gamma(1+a-b-c-d)}
 {\Gamma(1+a)\Gamma(1+a-b-c)\Gamma(1+a-b-d)
       \Gamma(1+a-c-d)}.
 \label{dr:eq:classical}
\end{align}
Here $\Rea(1+a-b-c-d)>0$ \cite{dlmf-dougall}. Substitute
$a=2\kappa$ and $d=d_0-u$. The ratio
$(1+a/2)_n/(a/2)_n=(n+\kappa)/\kappa$ produces the factor $x_n$.
The normalization in \eqref{dr:eq:weight}, followed by
$\Gamma(1+2\kappa)=2\kappa\Gamma(2\kappa)$, leaves exactly the factor
$1/2$ in \eqref{dr:eq:gamma-side}.

The sum of the lower hypergeometric parameters minus the sum of the upper
ones is $2u$. Thus the usual hypergeometric excess is $2u$, even though the
Gamma factor that detects the integral resonances is $\Gamma(u)$.
This distinction controls every factor of two below.

\section{Centered parity and normally convergent subtraction}
\label{dr:sec:subtraction}

\subsection{The coefficient polynomials}

Set
\begin{equation}
 r_1=\kappa,\quad r_2=b-\kappa,\quad r_3=c-\kappa,
 \quad r_4=1+\kappa-b-c-u.
 \label{dr:eq:r}
\end{equation}
In the coordinate $x=x_n$ the summand becomes
\begin{equation}
 W_n(u)=x\prod_{i=1}^4
           \frac{\Gamma(x+r_i)}{\Gamma(x+1-r_i)}.
 \label{dr:eq:centered}
\end{equation}
Define polynomials $\ell_j$ and $A_j$ by
\begin{align}
 \ell_j(u)&=-\frac1{j(2j+1)}\sum_{i=1}^4B_{2j+1}(r_i),
       &&j\geq1,\label{dr:eq:ell}\\
 \sum_{j\geq0}A_j(u)y^j
   &=\exp\!\left(\sum_{j\geq1}\ell_j(u)y^j\right).
       \label{dr:eq:A}
\end{align}
The second identity is formal in $y$; convergence of the infinite
asymptotic expansion is not asserted. Its finite recursion is
\begin{equation}
 A_0=1,\qquad A_j=\frac1j\sum_{r=1}^j r\ell_r A_{j-r}.
 \label{dr:eq:recursion}
\end{equation}
For example,
\begin{equation}
 A_1=\ell_1,\qquad A_2=\ell_2+\frac{\ell_1^2}{2},\qquad
 A_3=\ell_3+\ell_1\ell_2+\frac{\ell_1^3}{6}.
 \label{dr:eq:first-A}
\end{equation}
These expressions are polynomials with rational coefficients in
$\kappa,b,c,u$.

\begin{lemma}[Even inverse-power expansion]
\label{dr:lem:parity}
For every fixed positive integer $K$,
\begin{equation}
 W_n(u)=x_n^{-1-2u}
 \left(\sum_{j=0}^{K-1}A_j(u)x_n^{-2j}
             +O(x_n^{-2K})\right).
 \label{dr:eq:asymptotic}
\end{equation}
The remainder is locally uniform in the admissible parameters and in
$u$ with $\Rea u<d_0$. After $m$ spectral derivatives, the error is
$O(x_n^{-1-2\Rea u-2K}(1+\log^m x_n))$ on such compact sets.
\end{lemma}

\begin{proof}
The classical logarithmic Gamma-ratio expansion is
\begin{align*}
 \log\frac{\Gamma(x+r)}{\Gamma(x+1-r)}
 \sim{}&(2r-1)\log x\\
 &+\sum_{k\geq1}\frac{(-1)^{k+1}
          \bigl(B_{k+1}(r)-B_{k+1}(1-r)\bigr)}{k(k+1)x^k}.
\end{align*}
It follows from the shifted Stirling expansion, uniformly for $r$ in
compact sets \cite{dlmf-gamma}. Since
$B_m(1-r)=(-1)^mB_m(r)$, the coefficient is zero when $k$ is odd. For
$k=2j$ it equals $-B_{2j+1}(r)/(j(2j+1))$.
Furthermore $\sum_i r_i=1-u$, so the prefactor $x$ in
\eqref{dr:eq:centered} gives total exponent
$1+\sum_i(2r_i-1)=-1-2u$. Summation and formal exponentiation give
\eqref{dr:eq:asymptotic} and \eqref{dr:eq:recursion}.

For the derivative statement, first multiply the difference in
\eqref{dr:eq:asymptotic} by $x_n^{1+2u}$. Its analytic remainder is
$O(x_n^{-2K})$ uniformly in a slightly larger parameter neighborhood.
Cauchy's estimates give the same order for its fixed spectral
derivatives. Differentiating the removed factor $x_n^{-1-2u}$ then
contributes only powers of $-2\log x_n$. The finitely many initial
summands do not affect the argument.
\end{proof}

\begin{definition}[Specified subtraction]
For $K\geq1$, define
\begin{equation}
 \Rcal_K(u)=\sum_{n=0}^{\infty}
 \left[W_n(u)-\sum_{j=0}^{K-1}
           A_j(u)x_n^{-1-2u-2j}\right].
 \label{dr:eq:R-definition}
\end{equation}
The bracket is a single summand; its components need not be separately
summable.
\end{definition}

\begin{theorem}[Holomorphic subtraction identity]
\label{dr:thm:subtraction}
The series \eqref{dr:eq:R-definition} converges absolutely and normally
on
\begin{equation}
 -K<\Rea u<d_0.
 \label{dr:eq:strip}
\end{equation}
Every fixed spectral derivative can be taken term by term. On that strip,
with removable values understood,
\begin{equation}
 \Rcal_K(u)=G(u)-\sum_{j=0}^{K-1}
                      A_j(u)\zeta(1+2u+2j,\kappa).
 \label{dr:eq:subtraction}
\end{equation}
\end{theorem}

\begin{proof}
On a compact subset of \eqref{dr:eq:strip}, Lemma~\ref{dr:lem:parity}
majorizes the bracket and each fixed derivative by a constant times
$n^{-1-\epsilon}(1+\log^m n)$ for some $\epsilon>0$. This proves normal
convergence and holomorphy. On the nonempty substrip $0<\Rea u<d_0$,
all terms can be summed separately, and \eqref{dr:eq:dougall} gives
\eqref{dr:eq:subtraction}. The identity theorem extends it throughout
\eqref{dr:eq:strip}. Any poles on its displayed right side must be
removable because the left side is holomorphic.
\end{proof}

This theorem constructs an ordinary convergent function, not an
unspecified finite part. In particular, at $u=-N$ the minimal generic
choice $K=N+1$ leaves summands of size $O(n^{-3})$, and its $m$th
spectral derivative has summands $O(n^{-3}(1+\log^m n))$. The estimates
serve only to justify the exact identities.

\section{The residue polynomial and complete constant-term collapse}
\label{dr:sec:collapse}

\begin{theorem}[Explicit residue and coefficient differential law]
\label{dr:thm:coefficient-law}
Let $N\geq0$ and define
\begin{equation}
 \rho_N=\frac{(-1)^N}{N!}(d_0)_N(1-b)_N(1-c)_N.
 \label{dr:eq:rho}
\end{equation}
Then
\begin{equation}
 A_N(-N)=\rho_N.
 \label{dr:eq:residue}
\end{equation}
With $D=\partial_\kappa+\partial_b+\partial_c$, acting on the
polynomials with $d_0=1+2\kappa-b-c$, one also has
\begin{equation}
 (\partial_u-D)A_N(u)=
      \sum_{r=1}^N\frac{B_{2r}(\kappa)}r A_{N-r}(u).
 \label{dr:eq:law}
\end{equation}
The sum is empty when $N=0$.
\end{theorem}

\begin{proof}
The residue of $G(u)$ at $u=-N$ is $\rho_N/2$, since
\[
 \frac{\Gamma(b)}{\Gamma(b-N)}=(-1)^N(1-b)_N,
 \qquad \frac{\Gamma(d_0+N)}{\Gamma(d_0)}=(d_0)_N.
\]
The Gamma quotient identities hold by polynomial continuation at zeros.
In \eqref{dr:eq:subtraction} with $K=N+1$, only the $j=N$ zeta
counterterm can have a pole there. Its residue as a function of $u$ is
$A_N(-N)/2$. Holomorphy proves \eqref{dr:eq:residue}. Since both sides
are polynomials in the fixed parameters, the identity is a polynomial
identity, not merely a generic numerical equality.

For the differential law observe that $D d_0=0$. The operator
$E=\partial_u-D$ acts on the four variables in \eqref{dr:eq:r} by
\[
 Er_1=-1,\qquad Er_2=Er_3=Er_4=0.
\]
Using $B_j'=jB_{j-1}$ in \eqref{dr:eq:ell} gives
$E\ell_r=B_{2r}(\kappa)/r$. Apply $E$ to the formal exponential
\eqref{dr:eq:A}, multiply its two formal series, and compare the
coefficient of $y^N$. This proves \eqref{dr:eq:law} at every order.
\end{proof}

\begin{theorem}[Collapsed value at every integral resonance]
\label{dr:thm:collapse}
For every integer $N\geq0$, under \eqref{dr:eq:domain},
\begin{align}
 \Rcal_{N+1}(-N)=\frac{\rho_N}{2}
 \bigl[\psi(N+1)-\psi(d_0+N)-\psi(b)-\psi(c)
                  +2\psi(\kappa)\bigr].
 \label{dr:eq:collapse}
\end{align}
This formula remains valid when $b$ or $c$ belongs to $\{1,\ldots,N\}$;
its right side is then zero, without any singular digamma factor.
\end{theorem}

\begin{proof}
First assume $b-N$ and $c-N$ are not nonpositive integers. Near $u=-N$,
write $u=-N+t$ and expand
\begin{equation}
 G(-N+t)=\frac{\rho_N}{2t}
 +\frac{\rho_N}{2}
 \bigl[\psi(N+1)-\psi(d_0+N)-\psi(b-N)-\psi(c-N)\bigr]
 +O(t).
 \label{dr:eq:G-constant}
\end{equation}
The polar counterterm has constant coefficient
\[
 A_N(-N+t)\zeta(1+2t,\kappa)
 =\frac{\rho_N}{2t}+\frac12A_N'(-N)-\rho_N\psi(\kappa)+O(t).
\]
For the remaining terms, the Bernoulli special values give
\[
 \sum_{j=0}^{N-1}A_j(-N)\zeta(1-2(N-j),\kappa)
 =-\frac12\sum_{r=1}^N\frac{B_{2r}(\kappa)}r A_{N-r}(-N).
\]
Substitute these three expressions into \eqref{dr:eq:subtraction}.
Equation~\eqref{dr:eq:law} cancels the entire Bernoulli sum against
$A_N'(-N)-D\rho_N$. The surviving value is the constant coefficient
in \eqref{dr:eq:G-constant}, minus $D\rho_N/2$, plus
$\rho_N\psi(\kappa)$.

Since $D d_0=0$, logarithmic differentiation of the two remaining rising
factorials gives
\[
 D\rho_N=\rho_N\sum_{r=1}^N
           \left(\frac1{b-r}+\frac1{c-r}\right).
\]
Finally,
$\psi(b)-\psi(b-N)=\sum_{r=1}^N(b-r)^{-1}$, and similarly for $c$.
This proves \eqref{dr:eq:collapse} generically. Both the normally
convergent left side and the final right side are continuous, indeed
locally analytic, throughout \eqref{dr:eq:domain}. The final right side
contains digamma functions only at positive arguments. Taking limits
therefore proves the asserted degenerate cases as well.
\end{proof}

The zeroth resonance is particularly transparent:
\begin{equation}
 \sum_{n=0}^{\infty}\left[W_n(0)-\frac1{n+\kappa}\right]
 =\psi(\kappa)-\frac12\bigl[\gamma+\psi(d_0)+\psi(b)+\psi(c)\bigr].
 \label{dr:eq:Nzero}
\end{equation}
At higher resonances the raw subtracted expression contains several
Bernoulli polynomials and derivatives of asymptotic coefficients.
Theorem~\ref{dr:thm:collapse} says that all of them disappear from its
final value. The differential law, rather than a finite check at selected
$N$, explains this cancellation.

\section{Every spectral jet, including reciprocal-Gamma zeros}
\label{dr:sec:jets}

\subsection{A holomorphic Gamma factor}

For $N\geq0$ define the analytic germ
\begin{equation}
 Q_N(t)=\frac{(-1)^N}{N!}\Gamma(1+t)
       \prod_{r=1}^N(1-t/r)^{-1}P(-N+t).
 \label{dr:eq:Q}
\end{equation}
The empty product is one. The elementary Gamma recurrence gives
\begin{equation}
 G(-N+t)=\frac{Q_N(t)}{2t},\qquad Q_N(0)=\rho_N.
 \label{dr:eq:Q-identity}
\end{equation}
The factors $1/\Gamma(b-N+t)$ and $1/\Gamma(c-N+t)$ in $P$ remain
entire factors. Formula~\eqref{dr:eq:Q} therefore defines all jets even
when $\rho_N=0$.

Put $a_j(t)=A_j(-N+t)$. The subtraction theorem and
\eqref{dr:eq:stieltjes} yield the following analytic identity of germs:
\begin{align}
 \Rcal_{N+1}(-N+t)={}&\frac{Q_N(t)-a_N(t)}{2t}
 -a_N(t)\sum_{k\geq0}\frac{(-2)^k}{k!}\gamma_k(\kappa)t^k\notag\\
 &-\sum_{j=0}^{N-1}a_j(t)
              \zeta(1-2(N-j)+2t,\kappa).
 \label{dr:eq:jet-generator}
\end{align}
The first quotient is holomorphic because its numerator vanishes at zero.
There is no division by $\rho_N$.

\begin{theorem}[Finite closure of all integral-resonance jets]
\label{dr:thm:jets}
Write $q_r=[t^r]Q_N(t)$ and $a_{j,r}=[t^r]A_j(-N+t)$.
For every integer $m\geq0$,
\begin{align}
 \frac{\Rcal_{N+1}^{(m)}(-N)}{m!}
 ={}&\frac{q_{m+1}-a_{N,m+1}}2
 -\sum_{r=0}^m a_{N,r}
                 \frac{(-2)^{m-r}\gamma_{m-r}(\kappa)}{(m-r)!}
                 \notag\\
 &-\sum_{j=0}^{N-1}\sum_{r=0}^m a_{j,r}
    \frac{2^{m-r}\zeta^{(m-r)}(1-2(N-j),\kappa)}{(m-r)!}.
 \label{dr:eq:all-jets}
\end{align}
Thus the zeta spectral arguments outside the Stieltjes part belong to
$\{-1,-3,\ldots,1-2N\}$. Stieltjes orders are at most $m$, and the
Gamma germ is required only through Taylor degree $m+1$.
\end{theorem}

\begin{proof}
Take the coefficient of $t^m$ in \eqref{dr:eq:jet-generator}. All germs
are holomorphic near zero, so ordinary Cauchy products apply. The
termwise differentiated convergent series represents the same coefficient
by Theorem~\ref{dr:thm:subtraction}. The finite formula remains valid at
reciprocal-Gamma zeros because \eqref{dr:eq:Q} is their analytic Taylor
expansion, not a logarithmic expansion divided by a zero value.
\end{proof}

\subsection{Bell-polynomial formulas away from zeros}

We normalize complete exponential Bell polynomials by
\begin{equation}
 \exp\left(\sum_{r\geq1}h_r\frac{t^r}{r!}\right)
       =\sum_{r\geq0}Y_r(h_1,\ldots,h_r)\frac{t^r}{r!},
 \qquad Y_0=1.
 \label{dr:eq:Bell}
\end{equation}
When $b-N$ and $c-N$ avoid nonpositive integers, the logarithmic
coefficients of $Q_N$ are
\begin{align}
 h_{N,r}={}&\psi^{(r-1)}(1)+(r-1)!H_N^{(r)}
 +(-1)^r\psi^{(r-1)}(d_0+N)\notag\\
 &-\psi^{(r-1)}(b-N)-\psi^{(r-1)}(c-N),
 \label{dr:eq:hN}
\end{align}
where $H_N^{(r)}=\sum_{j=1}^N j^{-r}$. Consequently
\begin{equation}
 q_r=\frac{\rho_N}{r!}Y_r(h_{N,1},\ldots,h_{N,r}).
 \label{dr:eq:q-Bell}
\end{equation}
This convenient form is not the definition at a zero of $\rho_N$.
The safe computational alternative is to shift every problematic
reciprocal Gamma factor to a positive argument:
\begin{equation}
 \frac1{\Gamma(w+t)}=
 \frac{\prod_{j=0}^{L-1}(w+j+t)}{\Gamma(w+L+t)},
 \qquad w+L>0.
 \label{dr:eq:rgamma-safe}
\end{equation}
Expand the polynomial numerator and the nonzero Gamma germ separately.
The code implements precisely this operation.

\subsection{The zero-excess family in closed Bell form}

At $N=0$, $A_0=1$ and $Q_0(0)=1$. Set
\begin{equation}
 h_r=\psi^{(r-1)}(1)+(-1)^r\psi^{(r-1)}(d_0)
      -\psi^{(r-1)}(b)-\psi^{(r-1)}(c).
 \label{dr:eq:hzero}
\end{equation}
The general formula simplifies to
\begin{equation}
 \Rcal_1^{(m)}(0)=
 \frac{Y_{m+1}(h_1,\ldots,h_{m+1})}{2(m+1)}
                 -(-2)^m\gamma_m(\kappa).
 \label{dr:eq:zero-jets}
\end{equation}
For instance,
\begin{align}
 \Rcal_1(0)&=h_1/2-\gamma_0(\kappa),\notag\\
 \Rcal_1'(0)&=(h_1^2+h_2)/4+2\gamma_1(\kappa),\notag\\
 \Rcal_1''(0)&=(h_1^3+3h_1h_2+h_3)/6-4\gamma_2(\kappa).
 \label{dr:eq:first-jets}
\end{align}
No integer-relation search is involved in these identities.

\section{Half-integral values and exact polynomial degeneracies}
\label{dr:sec:half}

\subsection{The interlaced even-argument tower}

\begin{theorem}[Half-integral spectral values]
\label{dr:thm:half}
For $N\geq0$, put $u_*=-N-\tfrac12$. Then
\begin{equation}
 \Rcal_{N+1}(u_*)=G(u_*)+
 \sum_{j=0}^N A_j(u_*)
       \frac{B_{2(N-j)+1}(\kappa)}{2(N-j)+1}.
 \label{dr:eq:half}
\end{equation}
Every fixed spectral derivative at $u_*$ is a finite expression in Gamma
jets and Hurwitz-zeta derivatives at $0,-2,\ldots,-2N$; no Stieltjes
Laurent coefficient is needed.
\end{theorem}

\begin{proof}
At $u_*$ the counterterm spectral argument is $2(j-N)$, never one.
Also $\Gamma(u_*)$ is finite. Substitute directly in
\eqref{dr:eq:subtraction} and use
$\zeta(-m,a)=-B_{m+1}(a)/(m+1)$. Ordinary differentiation of the same
holomorphic identity gives the jet statement. Here the minimal subtraction
has terms $O(n^{-2})$, and the derivatives have the corresponding
logarithmic factors.
\end{proof}

The first case is
\begin{equation}
 \Rcal_1(-\tfrac12)=\kappa-\tfrac12-
 \sqrt\pi\frac{\Gamma(b)\Gamma(c)\Gamma(d_0+\tfrac12)}
 {\Gamma(d_0)\Gamma(b-\tfrac12)\Gamma(c-\tfrac12)}.
 \label{dr:eq:first-half}
\end{equation}
The different parity of the two zeta towers is a consequence of centered
very-well-poised symmetry, not a rule that applies to arbitrary Gamma
ratios.

\subsection{The symmetric half-parameter line}

Take $\kappa=b=c=\tfrac12$, so $d_0=1$. The summand simplifies to
\begin{equation}
 W_n(u)=\frac{\Gamma(n+1-u)}{(n+\tfrac12)\Gamma(n+1+u)}.
 \label{dr:eq:half-weight}
\end{equation}
At negative integer and negative half-integer parameters, the Gamma
recurrence gives two exact finite products. Writing $x=n+\tfrac12$,
\begin{align}
 W_n(-N)&=\frac1x\prod_{r=1}^N\left(x^2-(r-\tfrac12)^2\right),
        \label{dr:eq:integer-product}\\
 W_n(-N-\tfrac12)&=\prod_{r=1}^N(x^2-r^2).
        \label{dr:eq:half-product}
\end{align}
These include any initial zero summands, by reciprocal-Gamma continuation.
Their asymptotic expansions terminate identically. In particular,
\begin{equation}
 \Rcal_{N+1}(-N)=\Rcal_{N+1}(-N-\tfrac12)=0
 \quad\text{on this parameter line}.
 \label{dr:eq:zero-remainders}
\end{equation}
For the integer case,
\begin{equation}
 A_j(-N)=(-1)^j e_j\bigl((\tfrac12)^2,(\tfrac32)^2,
                         \ldots,(N-\tfrac12)^2\bigr),
 \label{dr:eq:elementary-coeff}
\end{equation}
with $A_j(-N)=0$ for $j>N$. Its final coefficient is
$\rho_N=(-1)^N(\tfrac12)_N^2$.

\begin{proposition}[An additional exact polynomial locus]
\label{dr:prop:integer-b}
If $N\geq1$ and $b=m\in\{1,\ldots,N\}$, then, under
\eqref{dr:eq:domain},
\begin{equation}
 W_n(-N)=x_n
 \prod_{r=1}^{m-1}\bigl(x_n^2-(\kappa-r)^2\bigr)
 \prod_{r=1}^{N-m}\bigl(x_n^2-(\kappa-c+r)^2\bigr).
 \label{dr:eq:integer-b}
\end{equation}
Thus its centered expansion ends at $j=N-1$ and
$\Rcal_{N+1}(-N)=0$ term by term. There is a symmetric assertion with
$b$ and $c$ exchanged.
\end{proposition}

\begin{proof}
In \eqref{dr:eq:weight}, combine the first pair of relevant Gamma
quotients as
\[
 \frac{\Gamma(n+m)}{\Gamma(n+1)}
 \frac{\Gamma(n+2\kappa)}{\Gamma(n+1+2\kappa-m)}
 =\prod_{r=1}^{m-1}(n+r)(n+2\kappa-r).
\]
The other pair, with $q=N-m$, is
\[
 \frac{\Gamma(n+c)}{\Gamma(n+c-q)}
 \frac{\Gamma(n+1+2\kappa-c+q)}{\Gamma(n+1+2\kappa-c)}
 =\prod_{r=1}^{q}(n+c-r)(n+2\kappa-c+r).
\]
Center each pair at $x_n$. Polynomial continuation handles all zero
reciprocal Gamma values. This proves \eqref{dr:eq:integer-b} and the
remaining conclusions.
\end{proof}

A specialized polynomial identity cannot be differentiated in $u$ after
$u$ has been fixed. The transverse spectral derivatives contain information
lost by that specialization. The next section gives a concrete example.

\section{Exact harmonic-number identities}
\label{dr:sec:harmonic}

\subsection{The general all-order harmonic form}

For $r\geq1$, let
\begin{equation}
 g_{n,r}=(-1)^r\psi^{(r-1)}(n+d_0)
                   -\psi^{(r-1)}(n+b+c).
 \label{dr:eq:g}
\end{equation}
Differentiating the logarithm of \eqref{dr:eq:weight} at zero gives
\begin{equation}
 \left.\partial_u^mW_n(u)\right|_{u=0}
   =W_n(0)Y_m(g_{n,1},\ldots,g_{n,m}).
 \label{dr:eq:W-Bell}
\end{equation}
Hence \eqref{dr:eq:zero-jets} is the ordinary convergent identity
\begin{align}
 &\sum_{n=0}^{\infty}
 \left[W_n(0)Y_m(g_{n,1},\ldots,g_{n,m})
               -\frac{(-2\log x_n)^m}{x_n}\right]\notag\\
 &\hspace{20mm}=
 \frac{Y_{m+1}(h_1,\ldots,h_{m+1})}{2(m+1)}
                  -(-2)^m\gamma_m(\kappa).
 \label{dr:eq:general-harmonic}
\end{align}
The polygamma differences are generalized harmonic numbers. With
$H_n^{(r)}(a)=\sum_{j=0}^{n-1}(j+a)^{-r}$,
\begin{equation}
 \psi^{(r-1)}(n+a)=\psi^{(r-1)}(a)
                  +(-1)^{r-1}(r-1)!H_n^{(r)}(a).
 \label{dr:eq:harmonic-conversion}
\end{equation}
Thus \eqref{dr:eq:general-harmonic} is an explicit all-order family of
Gamma-weighted harmonic identities. The compact normal convergence proof
already justifies every coefficient interchange.

\subsection{A half-parameter family without Gamma weights}

On the symmetric line of \eqref{dr:eq:half-weight}, $W_n(0)=1/x_n$.
Put $L=\log2$ and $A=\gamma+2L$. The right-side Bell data become
\begin{equation}
 h_1=2A,\qquad
 h_r=\begin{cases}
 -2(2^r-2)(r-1)!\zeta(r),&r\geq2\text{ even},\\
  2(2^r-1)(r-1)!\zeta(r),&r\geq3\text{ odd}.
 \end{cases}
 \label{dr:eq:half-h}
\end{equation}
The left-side data are
\begin{equation}
 g_{n,1}=-2(H_n-\gamma),\qquad
 g_{n,r}=\begin{cases}
 0,&r\geq2\text{ even},\\
 2(r-1)!\bigl(\zeta(r)-H_n^{(r)}\bigr),&r\geq3\text{ odd}.
 \end{cases}
 \label{dr:eq:half-g}
\end{equation}
Here $H_n=\sum_{j=1}^n j^{-1}$ and $H_0=0$. Substitution in
\eqref{dr:eq:general-harmonic} gives every order in ordinary harmonic
numbers and zeta tails, with no Gamma weights remaining.

At order one this yields
\begin{equation}
 \sum_{n=0}^{\infty}
 \frac{H_n-\gamma-\log(n+\tfrac12)}{n+\tfrac12}
 =\frac{\zeta(2)-\gamma^2-2L^2-2\gamma_1}{2}.
 \label{dr:eq:harmonic-one}
\end{equation}
At order two it gives
\begin{equation}
 \sum_{n=0}^{\infty}
 \frac{(H_n-\gamma)^2-\log^2(n+\tfrac12)}{n+\tfrac12}
 =\frac{A^3}{3}-A\zeta(2)+\frac76\zeta(3)-\gamma_2(\tfrac12).
 \label{dr:eq:harmonic-two}
\end{equation}
For reference, differentiation of
$\zeta(s,\tfrac12)=(2^s-1)\zeta(s)$ at its Laurent pole gives
\begin{align}
 \gamma_1(\tfrac12)&=\gamma_1-2\gamma L-L^2,\notag\\
 \gamma_2(\tfrac12)&=\gamma_2-4L\gamma_1+2\gamma L^2+\frac23L^3.
 \label{dr:eq:half-stieltjes}
\end{align}
The first summands in \eqref{dr:eq:harmonic-one} are $O(n^{-3})$;
those in \eqref{dr:eq:harmonic-two} are $O(n^{-3}\log n)$.

\subsection{A first-resonance identity with a revived head}

The same half-parameter family has
\begin{equation}
 A_1(u)=\frac{u(4u^2-1)}{12},\quad
 A_1(-1)=-\frac14,\quad A_1'(-1)=\frac{11}{12},\quad
 A_1''(-1)=-2.
 \label{dr:eq:A1-half}
\end{equation}
For $n\geq1$,
\begin{align}
 W_n(-1)&=\frac{n(n+1)}{n+\tfrac12},\notag\\
 W_n'(-1)&=\frac{n(n+1)}{n+\tfrac12}
       \left[-2(H_n-\gamma)+\frac1{n(n+1)}\right].
 \label{dr:eq:W1-half}
\end{align}
The initial term must be computed separately. In a local coordinate $t$,
\begin{equation}
 W_0(-1+t)=\frac{2\Gamma(2-t)}{\Gamma(t)}=2t+O(t^2).
 \label{dr:eq:revived-head}
\end{equation}
Thus $W_0(-1)=0$ but $W_0'(-1)=2$. After subtracting the derivative
of the two counterterms, the initial bracket is $1/6$.

\begin{corollary}[An explicit $\zeta'(-1)$ harmonic identity]
\label{dr:cor:first-resonance}
One has
\begin{align}
 &\frac16+\sum_{n=1}^{\infty}
 \left[\frac{2n(n+1)}{n+\tfrac12}
       \bigl(\log(n+\tfrac12)-H_n+\gamma\bigr)
       +\frac1{12(n+\tfrac12)}\right]\notag\\
 &\qquad=\zeta'(-1)+\frac{\zeta(2)}4-\frac{\gamma^2}4
        -\frac{\gamma_1}2-\frac{L^2}2+\frac\gamma{12}+\frac L4.
 \label{dr:eq:first-resonance-harmonic}
\end{align}
\end{corollary}

\begin{proof}
The bracket is $\partial_u\Rcal_2(u)|_{u=-1}$ by
\eqref{dr:eq:A1-half}--\eqref{dr:eq:revived-head}. For this specialization,
\begin{equation}
 Q_1(t)=-\frac{\pi\Gamma(1+t)\Gamma(1-t)}{\Gamma(-\tfrac12+t)^2}.
 \label{dr:eq:Q1-half}
\end{equation}
Insert its Taylor coefficients and \eqref{dr:eq:A1-half} into
\eqref{dr:eq:all-jets} with $N=m=1$. Direct simplification gives
\[
 -\frac{A^2}{4}+\frac A{12}+\frac{\zeta(2)}4
 -\frac{\gamma_1(\tfrac12)}2-2\zeta'(-1,\tfrac12).
\]
Use \eqref{dr:eq:half-stieltjes} and
$\zeta'(-1,\tfrac12)=-\zeta'(-1)/2-L/24$. This is precisely the
right side of \eqref{dr:eq:first-resonance-harmonic}. Normal convergence
of the differentiated subtraction proves ordinary convergence of the
printed series.
\end{proof}

The values of the left sides of \eqref{dr:eq:harmonic-one},
\eqref{dr:eq:harmonic-two}, and \eqref{dr:eq:first-resonance-harmonic}
begin, respectively,
\[
 0.24824090308572918127,\qquad
 -0.27295567580201109333,\qquad
 0.18008740839086106049.
\]
These decimals are diagnostic illustrations of the proved formulas.

\section{Continuous shifts and exact Euler-primitive endpoints}
\label{dr:sec:harmonic-primitives}

For $a>0$ and integer $p\geq1$, let
\begin{equation}
 S_p(a)=\sum_{n=0}^{\infty}
       \frac{H_n-\gamma-\log(n+a)}{(n+a)^p}.
 \label{dr:eq:Sp}
\end{equation}
The numerator is $O(n^{-1})$, locally uniformly in $a$, so these are
ordinary absolutely convergent sums even at $p=1$.

\begin{theorem}[A continuous-shift harmonic closure]
\label{dr:thm:Sp}
For every $a>0$,
\begin{equation}
 S_1(a)=\frac{\zeta(2,a)-\psi(a)^2}{2}-\gamma_1(a)-\zeta(2).
 \label{dr:eq:S1-general}
\end{equation}
For every integer $p\geq2$,
\begin{align}
 S_p(a)={}&\psi(a)\zeta(p,a)+\frac p2\zeta(p+1,a)
              +\zeta'(p,a)\notag\\
 &-\frac12\sum_{j=2}^{p-1}\zeta(j,a)\zeta(p+1-j,a).
 \label{dr:eq:Sp-general}
\end{align}
An empty sum is zero.
\end{theorem}

\begin{proof}
Begin with a classical beta differentiation. The expansion at $v=0$ is
\begin{equation}
 B(a,v)=\frac1v-(\gamma+\psi(a))
 +\frac v2\bigl((\gamma+\psi(a))^2+\zeta(2)-\psi'(a)\bigr)
 +O(v^2).
 \label{dr:eq:beta}
\end{equation}
For $v>0$, differentiation under Euler's beta integral gives
$\partial_a\partial_vB(a,v)$. At $v=0$ this mixed derivative remains
an ordinary convergent integral: its integrand is
$t^{a-1}\log t\log(1-t)/(1-t)$, which is integrable at both endpoints.
Using
$\sum_{n\geq0}H_nt^n=-\log(1-t)/(1-t)$, it follows that
\begin{equation}
 \sum_{n\geq0}\frac{H_n}{(n+a)^2}
       =(\psi(a)+\gamma)\psi'(a)-\frac12\psi''(a).
 \label{dr:eq:Hp2}
\end{equation}
The integral interchange can also be justified by monotone convergence,
since its expanded summands are nonnegative for $a>0$.

Differentiate \eqref{dr:eq:Hp2} $p-2$ times in $a$, and multiply by
$(-1)^{p-2}/(p-1)!$. The series derivatives converge normally on compact
subintervals of $a>0$. Substituting
$\psi^{(r)}(a)=(-1)^{r+1}r!\zeta(r+1,a)$ gives
\begin{equation}
 \sum_{n\geq0}\frac{H_n}{(n+a)^p}
 = (\psi(a)+\gamma)\zeta(p,a)+\frac p2\zeta(p+1,a)
 -\frac12\sum_{j=2}^{p-1}\zeta(j,a)\zeta(p+1-j,a).
 \label{dr:eq:Hp}
\end{equation}
For the last sum, the Leibniz coefficient before symmetrization is
$(p-j)/(p-1)$; the coefficient paired under $j\mapsto p+1-j$ adds to
one. Subtract $\gamma\zeta(p,a)$ and use
$\sum\log(n+a)/(n+a)^p=-\zeta'(p,a)$ to obtain
\eqref{dr:eq:Sp-general}.

To prove \eqref{dr:eq:S1-general}, termwise differentiation gives
$S_1'(a)=-\zeta(2,a)-S_2(a)$. Differentiating
\eqref{dr:eq:stieltjes} in $a$ and using
$\partial_a\zeta(s,a)=-s\zeta(s+1,a)$ gives
\begin{equation}
 \gamma_1'(a)=\zeta(2,a)+\zeta'(2,a).
 \label{dr:eq:gamma1-derivative}
\end{equation}
Consequently the derivative of
$\tfrac12(\zeta(2,a)-\psi(a)^2)-\gamma_1(a)$ is also
$-\zeta(2,a)-S_2(a)$. Their difference is constant on $a>0$.
Evaluate it at $a=1/2$ using \eqref{dr:eq:harmonic-one} and
\eqref{dr:eq:half-stieltjes}; the constant is $-\zeta(2)$.
\end{proof}

The beta-differentiation method and the underlying zeta identities are
classical. Their role here is to evaluate an entire primitive family of
the newly organized spectral identity, not to assert literature priority
for \eqref{dr:eq:Hp}.

At the half shift, \eqref{dr:eq:Sp-general} reduces entirely to Riemann
zeta values and derivatives:
\begin{align}
 S_p(\tfrac12)={}&(2^p-1)\zeta'(p)
 -\bigl((2^p-1)\gamma+(2^p-2)L\bigr)\zeta(p)\notag\\
 &+\frac p2(2^{p+1}-1)\zeta(p+1)\notag\\
 &-\frac12\sum_{j=2}^{p-1}(2^j-1)(2^{p+1-j}-1)
                       \zeta(j)\zeta(p+1-j),\qquad p\geq2.
 \label{dr:eq:Sp-half}
\end{align}
In particular,
\begin{align}
 S_2(\tfrac12)&=7\zeta(3)-(3\gamma+2L)\zeta(2)+3\zeta'(2),
       \label{dr:eq:S2-half}\\
 S_3(\tfrac12)&=\frac{45}{4}\zeta(4)
                   -(7\gamma+6L)\zeta(3)+7\zeta'(3).
       \label{dr:eq:S3-half}
\end{align}
Section~\ref{dr:sec:polylog} explains exactly why these are endpoint
values of successive normalized Euler antiderivatives. Notice the change
of coordinates: the unintegrated endpoint uses $\gamma_1$, but every
positive primitive order in this particular first-jet family uses
positive-integer zeta values and their first derivatives instead.

\section{Polylogarithmic subtraction and normalized primitives}
\label{dr:sec:polylog}

\subsection{Generating functions with their endpoint values retained}

For $|z|<1$ put
\begin{equation}
 F(u,z)=\sum_{n\geq0}W_n(u)z^n,
 \qquad
 \Phi(z,s,a)=\sum_{n\geq0}\frac{z^n}{(n+a)^s}.
 \label{dr:eq:F-Phi}
\end{equation}
The latter is Lerch's transcendent \cite{dlmf-lerch}. Define
\begin{equation}
 \Rcal_K(u,z)=F(u,z)-\sum_{j=0}^{K-1}
                    A_j(u)\Phi(z,1+2u+2j,\kappa).
 \label{dr:eq:Rz}
\end{equation}
Equivalently, it is the generating series of the brackets in
\eqref{dr:eq:R-definition}.

\begin{theorem}[Endpoint and spectral-jet lifting]
\label{dr:thm:endpoint}
On compact subsets of $-K<\Rea u<d_0$, the bracketed series defining
$\Rcal_K(u,z)$ and all its fixed spectral derivatives converge uniformly
for $|z|\leq1$. Consequently
\begin{equation}
 \partial_u^m\Rcal_K(u,1)=\Rcal_K^{(m)}(u).
 \label{dr:eq:endpoint}
\end{equation}
In particular, \eqref{dr:eq:collapse} and \eqref{dr:eq:all-jets} give
exact endpoint values of the corresponding subtracted hypergeometric--Lerch
identities.
\end{theorem}

\begin{proof}
Use the summable majorants in Theorem~\ref{dr:thm:subtraction} and
$|z^n|\leq1$. Uniform convergence gives continuity on the closed disk,
and the same argument applies to every fixed spectral derivative. This
is stronger than merely assigning an Abel value to a divergent series.
The unsubtracted $F$ and the individual Lerch terms need not have finite
endpoint values in this range.
\end{proof}

At generic parameters the hypergeometric form is
\begin{equation}
 F(u,z)=W_0(u)\,
 {}_5F_4\!\left(\begin{matrix}
 2\kappa,1+\kappa,b,c,d_0-u\\
 \kappa,1+2\kappa-b,1+2\kappa-c,b+c+u
 \end{matrix};z\right).
 \label{dr:eq:F-hyper}
\end{equation}
At a singular lower parameter this product can look like zero times an
undefined hypergeometric expression. The coefficient definition
\eqref{dr:eq:F-Phi}, with reciprocal-Gamma factors, is authoritative there.

\subsection{Rational centering converts the counterterms to polylogarithms}

Let $\kappa=p/q$, $1\leq p\leq q$, and
$\omega=\exp(2\pi i/q)$. For $0<z<1$, the root-of-unity filter gives
\begin{equation}
 \Phi(z,s,p/q)=q^{s-1}z^{-p/q}
       \sum_{\ell=0}^{q-1}\omega^{-p\ell}
                     \Li_s(\omega^\ell z^{1/q}).
 \label{dr:eq:root-filter}
\end{equation}
To prove it, expand each polylogarithm in its defining series. The finite
Fourier sum is $q$ when its summation index is congruent to $p$ modulo
$q$, and zero otherwise. The remaining terms are exactly those in
\eqref{dr:eq:F-Phi}. Analytic continuation in the punctured disk, using a
consistent root, gives the same single-valued analytic combination; the
apparent singularity at $z=0$ cancels. Rational $\kappa>1$ is reduced by
the finite shift formula
\begin{equation}
 \Phi(z,s,a+M)=z^{-M}\left(\Phi(z,s,a)
                  -\sum_{n=0}^{M-1}\frac{z^n}{(n+a)^s}\right).
 \label{dr:eq:lerch-shift}
\end{equation}

For spectral derivatives, differentiation of the factor $q^{s-1}$ must
be retained. If $s=s_0+2u$, then
\begin{align}
 \partial_u^m\Phi(z,s,p/q)
 ={}&2^m q^{s-1}z^{-p/q}\sum_{\ell=0}^{q-1}\omega^{-p\ell}
       (\partial_s+\log q)^m
                   \Li_s(\omega^\ell z^{1/q}).
 \label{dr:eq:root-jets}
\end{align}
Thus every rational-center counterterm and every counterterm jet has an
explicit finite polylogarithmic representation. This statement concerns
the subtraction terms; it does not assert a general reduction of the
remaining ${}_5F_4$ function to classical polylogarithms.

\subsection{Euler derivatives and genuine normalized antiderivatives}

Define
\begin{equation}
 \Theta_\kappa=z\frac{\mathrm d}{\mathrm dz}+\kappa,
 \qquad
 J_\kappa f(z)=z^{-\kappa}\int_0^z t^{\kappa-1}f(t)\dd t.
 \label{dr:eq:Euler}
\end{equation}
For a power series $f(z)=\sum f_nz^n$, this is the analytic operator
$J_\kappa f=\sum f_nz^n/(n+\kappa)$. It is the unique analytic inverse
of $\Theta_\kappa$ at zero, and
$\mathrm d(z^\kappa J_\kappa f)/\mathrm dz=z^{\kappa-1}f(z)$ is an
ordinary antiderivative identity. For integers $r\geq1$,
\begin{equation}
 J_\kappa^r f(z)=\frac{z^{-\kappa}}{(r-1)!}
     \int_0^z t^{\kappa-1}\bigl(\log(z/t)\bigr)^{r-1}f(t)\dd t.
 \label{dr:eq:iterated-primitive}
\end{equation}
Initially one may integrate on $0<z<1$; the power-series definition fixes
the analytic continuation without a branch ambiguity. It follows that
\begin{equation}
 \Theta_\kappa\Phi(z,s,\kappa)=\Phi(z,s-1,\kappa),
 \qquad J_\kappa^r\Phi(z,s,\kappa)=\Phi(z,s+r,\kappa).
 \label{dr:eq:Phi-primitive}
\end{equation}
Apply these operators coefficientwise to \eqref{dr:eq:Rz}. Spectral
derivatives commute with $J_\kappa^r$ on compact parameter sets, and
endpoint continuity for the bracketed primitives follows from absolute
convergence. Euler differentiation inside the disk is unrestricted;
any endpoint claim after differentiation additionally requires the
resulting coefficient series to be summable.

The first primitive of the generic hypergeometric side cancels the
very-well-poised pair:
\begin{align}
 J_\kappa F(u,z)={}&\frac{W_0(u)}\kappa\,
 {}_4F_3\!\left(\begin{matrix}2\kappa,b,c,d_0-u\\
 1+2\kappa-b,1+2\kappa-c,b+c+u\end{matrix};z\right).
 \label{dr:eq:4F3}
\end{align}
Again the original coefficient series gives its value at exceptional
lower parameters.

For the first spectral jet on the symmetric half-parameter line,
\begin{equation}
 -\frac12\partial_u\Rcal_1(0,z)
 =\sum_{n\geq0}\frac{H_n-\gamma-\log(n+\tfrac12)}{n+\tfrac12}z^n.
 \label{dr:eq:first-jet-generating}
\end{equation}
Hence all its normalized Euler-primitive endpoints are explicitly known:
\begin{equation}
 -\frac12\left[J_{1/2}^{\,r}\partial_u\Rcal_1(0,z)\right]_{z=1}
      =S_{r+1}(\tfrac12),\qquad r\geq0.
 \label{dr:eq:all-primitive-endpoints}
\end{equation}
Equations~\eqref{dr:eq:harmonic-one} and \eqref{dr:eq:Sp-half}
evaluate the entire family. This is a closed endpoint theorem for this
specific first jet, not an assertion that arbitrary primitive orders of
all higher Dougall jets have already been reduced.

\section{Antiderivatives of all collapsed resonance values}
\label{dr:sec:kappa-primitives}

The centering parameter gives a different antiderivative problem.
Here $b,c,N$ are fixed, while $\kappa$ varies within
\eqref{dr:eq:domain}. Theorem~\ref{dr:thm:collapse} makes this problem
finite because $\rho_N$ is a polynomial in $\kappa$.

\subsection{A normalized Hurwitz-zeta ladder}

For integers $j\geq0$ set
\begin{equation}
 \Pcal_j(x)=\frac{\zeta'(-j,x)+H_j\zeta(-j,x)}{j!},
 \qquad H_0=0.
 \label{dr:eq:Pj}
\end{equation}
Then
\begin{equation}
 \Pcal_0'(x)=\psi(x),\qquad
 \Pcal_j'(x)=\Pcal_{j-1}(x)\quad(j\geq1).
 \label{dr:eq:P-ladder}
\end{equation}
Indeed, Lerch's classical identity
$\zeta'(0,x)=\log\Gamma(x)-\tfrac12\log(2\pi)$ proves the first
relation. For the second, differentiate in the parameter to obtain
\[
 \partial_x\zeta'(-j,x)=j\zeta'(1-j,x)-\zeta(1-j,x),
 \qquad \partial_x\zeta(-j,x)=j\zeta(1-j,x),
\]
and use $H_j-1/j=H_{j-1}$. These normalized negative-polygamma
coordinates are classical in substance \cite{espinosa,dlmf-hurwitz}.

For a polynomial $p$, define
\begin{equation}
 \Ical[p](x)=\sum_{j=0}^{\deg p}(-1)^jp^{(j)}(x)\Pcal_j(x),
 \qquad \Ical[0]=0.
 \label{dr:eq:I}
\end{equation}
Adjacent terms telescope under differentiation, giving
\begin{equation}
 \frac{\mathrm d}{\mathrm dx}\Ical[p](x)=p(x)\psi(x).
 \label{dr:eq:I-derivative}
\end{equation}
This is a finite formula rather than an unevaluated repeated integral.

\begin{theorem}[A finite primitive at every resonance]
\label{dr:thm:primitive}
Let
\begin{align}
 C_N&=\frac{(-1)^N}{N!}(1-b)_N(1-c)_N,
 &p_N(\kappa)&=C_N(1+2\kappa-b-c)_N,\notag\\
 y&=1+2\kappa-b-c+N,
 &q_N(y)&=(y-N)_N,\notag\\
 B_N&=\psi(N+1)-\psi(b)-\psi(c).
 \label{dr:eq:primitive-data}
\end{align}
Choose any polynomial $P_N^{\rm poly}$ with derivative $p_N$.
An antiderivative, in $\kappa$, of $\Rcal_{N+1}(-N)$ is
\begin{equation}
 \frac{B_N}{2}P_N^{\rm poly}(\kappa)
   +\Ical[p_N](\kappa)-\frac{C_N}{4}\Ical[q_N](y).
 \label{dr:eq:all-kappa-primitives}
\end{equation}
The special-function coordinates required are only
$\zeta'(-j,\kappa)$, $\zeta'(-j,y)$ and Bernoulli polynomials,
$0\leq j\leq N$, together with the displayed constant digamma values.
\end{theorem}

\begin{proof}
Since $p_N(\kappa)=\rho_N=C_Nq_N(y)$ and $y'=2$, differentiating
\eqref{dr:eq:all-kappa-primitives} by \eqref{dr:eq:I-derivative} gives
\[
 \frac{B_N}{2}\rho_N+\rho_N\psi(\kappa)
                        -\frac{\rho_N}{2}\psi(y).
\]
This is exactly \eqref{dr:eq:collapse}. The Bernoulli assertion follows
from $\zeta(-j,x)=-B_{j+1}(x)/(j+1)$. If $C_N=0$, the displayed expression is constant and the identity
remains valid.
\end{proof}

At $N=0$, the formula reduces, up to a constant, to
\begin{equation}
 \log\Gamma(\kappa)-\frac14\log\Gamma(d_0)
       -\frac\kappa2\bigl(\gamma+\psi(b)+\psi(c)\bigr).
 \label{dr:eq:first-kappa-primitive}
\end{equation}
At $N=1$ only $\zeta'(0,\cdot)$ and $\zeta'(-1,\cdot)$ occur;
at each subsequent resonance precisely one further level is sufficient.
The theorem concerns zeroth spectral order at all resonances. Mixed
$\kappa$-primitives of arbitrary spectral jets are a separate reduction
problem, listed below.

\section{Audit, safeguards, and reproducible evidence}
\label{dr:sec:audit}

\subsection{What is, and is not, an erratum}

The inspected canonical statements already distinguish analytic proofs,
exact finite checks, and numerical residuals. No additional specific
mathematical error in those inspected statements was established by this
investigation. In particular, the outgoing package does not relabel a
classical identity as an error or silently change the status of an
unresolved Gaussian candidate.

There are, however, four essential corrections to a \emph{naive
transplant} of the Gauss construction. They are integration safeguards,
not accusations that the existing manuscript makes these mistakes.

\paragraph{The doubled excess.}
The counterterm is $\zeta(1+2u+2j,\kappa)$, with residue $1/2$ in
$u$. Replacing it by $\zeta(1+u+j,\kappa)$ would destroy the pole
cancellation and all derivative factors. The classical ${}_5F_4$ excess
is $2u$.

\paragraph{The centering and parity.}
Only the coordinate $x=n+\kappa$ has the complementary pairs used in
\eqref{dr:eq:centered}. An expansion in $n$ generally has odd inverse
powers. Omitting them without making the centering change is invalid.

\paragraph{Zeros are not zero jets.}
Equation~\eqref{dr:eq:revived-head} gives a concrete counterexample to
dropping a summand that vanishes at the expansion point. Its contribution
to the first derivative is $2$, and its complete subtracted contribution
is $1/6$. Likewise, the Bell formula \eqref{dr:eq:q-Bell} must not be
used as zero times divergent polygammas. Formula~\eqref{dr:eq:rgamma-safe}
is the replacement.

\paragraph{An asymptotic tail is not an interval certificate.}
The asymptotic coefficients are used in finite number in both the proof
and the numerical checks. Extending the asymptotic series to infinitely
many terms at fixed $n$ is not justified. Numerical agreement between
finite asymptotic-tail evaluations does not constitute a rigorous error
enclosure or establish arithmetic independence.

\subsection{Exact finite verification}

The delivered \file{code/verify_exact.py} completed \textbf{485 exact
assertions}. They include Bernoulli reflection through degree~24; generic
logarithmic differential identities through $j=6$; generic coefficient
and residue identities through $N=3$; rational residue and cancellation
checks at six parameter triples through $N=10$; and both symmetric
finite-product coefficient laws through $N=12$, including three
coefficients beyond the expected truncation at each level.

Additional checks cover the printed $A_1$ and initial-head data,
polynomial-weighted primitive telescoping, the factor-$1/2$ residue,
an independent truncated-exponential coefficient calculation, and the
product symmetrization in the Euler-primitive formula. All use exact
rational or symbolic polynomial arithmetic. Their scope is finite
regression verification; the all-order assertions are proved in the
preceding sections, and are not proof-assistant formalizations.

\subsection{Numerical diagnostics with independently constructed sides}

The numerical suites completed \textbf{293 comparisons at 70 working
decimal digits}. The central spectral tests use five parameter triples,
$N=0,1,2,3$, and Taylor orders $0,1,2,3$. Two series-side evaluations use
80 actual terms with coefficients through $A_{18}$ and 120 actual terms
with coefficients through $A_{22}$. Both are compared to the finite
Gamma--Stieltjes formula. The parameter triples include one and two
simultaneous reciprocal-Gamma zeros.

For a head length $M$, the series-side approximation is
\begin{align}
 &\sum_{n=0}^{M-1}\left[W_n(u)-\sum_{j<K}A_j(u)x_n^{-1-2u-2j}\right]
 +\sum_{j=K}^{J}A_j(u)\zeta(1+2u+2j,M+\kappa).
 \label{dr:eq:numeric-tail}
\end{align}
It uses the actual Gamma coefficient sequence, not the exact sum $G(u)$.
Its Taylor coefficients are computed by ordinary truncated-series
arithmetic. The three explicit harmonic series are also checked by a
separate digamma/Bernoulli tail construction, rather than merely calling
the same Dougall evaluator twice.

Other tests compare interior generating functions with independent
hypergeometric/Lerch evaluations; check rational root-of-unity filters
with order derivatives; differentiate the finite $\kappa$-primitives;
and evaluate the continuous-shift and half-parameter primitive endpoints.
The largest recorded scaled discrepancy was
\begin{equation}
 2.18454477481895\times10^{-48},\qquad
 \text{scaled discrepancy}=\frac{|\text{left}-\text{right}|}
                                    {\max(1,|\text{right}|)}.
 \label{dr:eq:max-error}
\end{equation}
For the 120-term, $A_{22}$ central spectral run it was
$1.54137402015147\times10^{-58}$. These are observations about the
executed computations, not certified counts of correct digits. Full
values, parameter choices, tolerances and versions are in
\file{results/numerical_checks.json}.

\subsection{Reproduction and delivery scope}

The tested environment used Python~3.13.5, SymPy~1.14.0 and
mpmath~1.3.0. From the extracted package directory, run
\begin{verbatim}
python -m pip install -r requirements.txt
python code/verify_exact.py
python code/verify_numeric.py --part 1
python code/verify_numeric.py --part 2
latexmk -pdf -interaction=nonstopmode -halt-on-error article.tex
\end{verbatim}
The numerical parts are separately executable to fit constrained tool
runtimes; together they perform the same assertions as the default full
run. The scripts regenerate their JSON reports. Reproduction changes
runtime metadata, so it changes the corresponding file hashes. No
network access is needed after installing the dependencies and TeX
packages.

The numerical implementation supports the real parameter domain tested
here and uses arbitrary precision controlled by mpmath. The article's
analytic continuation in $u$ is broader than the set of numerical samples.
The source audit, integration notes, manuscript-ready core section, result
files, and SHA-256 manifest accompany the TeX and compiled PDF.

\section{Further research questions}
\label{dr:sec:future}

\paragraph{1. Classify cancellation operators beyond this family.}
The operator $E=\partial_u-\partial_\kappa-\partial_b-\partial_c$
acts on just one centered Gamma pair and produces a Hurwitz Bernoulli
sequence. Determine necessary and sufficient conditions on a product of
complementary Gamma pairs for an analogous parameter vector field to
collapse all resonance constant terms. An exact classification of this
mechanism would be stronger than cataloguing isolated summations.

\paragraph{2. Compute multivariate resonant jets without path artifacts.}
The germ $Q_N$ handles every reciprocal-Gamma zero along one spectral
coordinate. Allow $b$ and $c$ to approach integral values simultaneously
with $u+N\to0$, and formulate a multivariate coefficient identity before
specializing any coordinate. The normally convergent subtracted series
supplies a natural holomorphic target. Determine which alternative
finite-part prescriptions agree with that target and which introduce
path-dependent constants.

\paragraph{3. Reduce higher harmonic-jet primitive endpoints.}
Equation~\eqref{dr:eq:all-primitive-endpoints} evaluates all primitive
orders for the first symmetric spectral jet. For the second and higher
jets, determine a closed generating system for the products of harmonic
numbers and odd zeta tails that arise. A useful first target is a finite
formula for every primitive of \eqref{dr:eq:harmonic-two}, with explicit
control of whether double zeta values or only products remain. Such a
reduction is not an independence statement about the resulting constants.

\paragraph{4. Integrate higher spectral jets in the center parameter.}
Theorem~\ref{dr:thm:primitive} solves the zeroth-order problem at every
$N$. Its higher-jet counterparts contain Stieltjes functions and
polygamma products. Develop normalized primitives through $(N,m)=(3,2)$,
then seek a finite all-order recursion. Keep additive constants and
parameter-derivative conventions explicit.

\paragraph{5. Compare local endpoint and centered counterterm bases.}
At rational $\kappa$, the present counterterms are finite combinations of
polylogarithms and their order derivatives. Construct their precise
connection with local expansions in powers and logarithms of $1-z$,
including constant terms at both integral and half-integral spectral
values. This would sharpen the comparison with classical hypergeometric
continuation and partial-sum formulas.

\paragraph{6. Classify exact truncation loci.}
The symmetric products \eqref{dr:eq:integer-product}--\eqref{dr:eq:half-product}
and Proposition~\ref{dr:prop:integer-b} are explicit loci where the centered
asymptotic expansion terminates. Determine all parameter loci for which
$W_n(u)$ is a Laurent polynomial in $n+\kappa$. Then classify which
transverse jets retain finite reductions and which produce nontrivial
zeta derivative towers.

\paragraph{7. Extract rational-level identities without nonreduction claims.}
Specialize the residue and jet formulas at small rational centers and
parameters, and convert the Hurwitz values into compatible cyclotomic or
Dirichlet-character coordinates. The goal is a short explicit identity
catalogue, with each coefficient obtained exactly. Failure of a smaller
numerical basis to reproduce a value is not a proof that the omitted
coordinates are necessary.

\paragraph{8. Formalize the algebraic and analytic layers separately.}
The formal exponential recursion and differential law are polynomial
identities suitable for a proof assistant. Their analytic interface is the
compact-normal convergence of the subtracted Gamma-ratio series.
Formalizing both would establish a machine-checked chain from classical
Dougall summation to the complete jet formulas. The present package
provides regression data, not that formalization.

\section{Conclusion}

The explicit Dougall target proposed in the incoming research programme
has a complete answer in the centered ${}_5F_4$ family studied here.
Complementary Gamma pairs enforce even inverse-power corrections; their
Bernoulli coefficients obey a differential law; that law removes all
lower counterterms from every integral-resonance constant. The residue
polynomial and full Taylor formula then organize Stieltjes constants,
negative-odd Hurwitz derivatives, and reciprocal-Gamma degeneracies in one
ordinary convergent calculus.

The interlaced half-integer points explain a distinct nonpositive-even
zeta tower. Harmonic specialization produces explicit identities with
$\gamma_1$, $\gamma_2$ and $\zeta'(-1)$, while beta differentiation
evaluates an entire family of Euler-primitive endpoints. Polynomial
weighting and the normalized Hurwitz ladder give antiderivatives at every
resonance in the center parameter. These are exact structural results;
the computational files make their finite coefficient and implementation
checks reproducible without changing their proof status.

\appendix
\section{Implementation details and coefficient conventions}
\label{dr:app:implementation}

The implementation stores ordinary Taylor coefficients. For arrays
$a=(a_0,\ldots,a_m)$ and $b=(b_0,\ldots,b_m)$, multiplication is truncated
convolution. If $\exp(L(t))=\sum e_jt^j$, then
\begin{equation}
 e_0=e^{L_0},\qquad
 e_j=\frac1j\sum_{r=1}^j rL_re_{j-r}.
 \label{dr:eq:expjet}
\end{equation}
This recursion is used both for Gamma germs and for the coefficient
generating function, with their different indices kept separate.

For $r_4=r_4(u_0)$ and $\ell\geq1$, the spectral Taylor coefficient
of the logarithmic asymptotic coefficient is especially simple:
\begin{equation}
 [t^\ell]\ell_j(u_0+t)=
 -\frac{(-1)^\ell}{j(2j+1)}
 \binom{2j+1}{\ell}B_{2j+1-\ell}(r_4),
 \quad 1\leq\ell\leq2j+1.
 \label{dr:eq:ell-jet}
\end{equation}
It vanishes beyond that degree. Thus all $A_j$ jets can be generated by
finite rational arithmetic at rational expansion points, without symbolic
Gamma manipulation.

For the harmonic checks, the independent expansion
\begin{equation}
 \psi(x+\tfrac12)-\log x
       \sim-\sum_{j\geq1}\frac{B_{2j}(\tfrac12)}{2j x^{2j}}
 \label{dr:eq:digamma-tail}
\end{equation}
is used at $x=n+1/2$. Squaring it together with the logarithmic cross term
gives the second harmonic tail. More generally,
\begin{equation}
 \psi(n+1)-\log(n+a)
 \sim\sum_{k\geq1}\frac{(-1)^{k+1}B_k(1-a)}{k(n+a)^k}.
 \label{dr:eq:shifted-digamma-tail}
\end{equation}
These are finite asymptotic approximations in the diagnostics, never
asserted to converge as infinite series at fixed $n$.

The numerical code asks mpmath for spectral derivatives through its
\texttt{zeta(..., derivative=r)} interface, rather than differentiating
a special-value shortcut that may have a different numerical code path.
The analytic formulas in the article remain the source of truth.
The manifest covers the delivered files; neither a checksum nor a passing
regression suite substitutes for review of the proofs.

\begin{thebibliography}{99}
\raggedright

\bibitem{dlmf-dougall}
NIST Digital Library of Mathematical Functions,
\emph{Generalized Hypergeometric Functions: Argument Unity},
\S16.4, especially Eq.~16.4.9 (Rogers--Dougall very-well-poised sum).
\url{https://dlmf.nist.gov/16.4}. Accessed 10 October 2026.

\bibitem{dlmf-gamma}
NIST Digital Library of Mathematical Functions,
\emph{Gamma Function: Asymptotic Expansions}, \S5.11.
\url{https://dlmf.nist.gov/5.11}. Accessed 10 October 2026.

\bibitem{dlmf-hurwitz}
NIST Digital Library of Mathematical Functions,
\emph{Hurwitz Zeta Function}, \S25.11.
\url{https://dlmf.nist.gov/25.11}. Accessed 10 October 2026.

\bibitem{dlmf-lerch}
NIST Digital Library of Mathematical Functions,
\emph{Lerch's Transcendent}, \S25.14.
\url{https://dlmf.nist.gov/25.14}. Accessed 10 October 2026.

\bibitem{buhring}
W.~B\"uhring, \emph{Partial sums of hypergeometric series of unit argument},
Proceedings of the American Mathematical Society \textbf{132} (2004),
407--415. Preprint arXiv:math/0311126 (2003).
\url{https://arxiv.org/abs/math/0311126}.

\bibitem{chufu}
W.~Chu and A.~M.~Fu, \emph{Dougall--Dixon formula and harmonic number
identities}, The Ramanujan Journal \textbf{18} (2009), 11--31.
DOI: \href{https://doi.org/10.1007/s11139-007-9044-6}{10.1007/s11139-007-9044-6}.

\bibitem{nicholson}
M.~Nicholson, \emph{Quadratic Transformations of Hypergeometric Function
and Series with Harmonic Numbers}, arXiv:1801.02428v3 (2019).
\url{https://arxiv.org/abs/1801.02428}.

\bibitem{espinosa}
O.~R.~Espinosa and V.~H.~Moll,
\emph{On some integrals involving the Hurwitz zeta function: part~2},
The Ramanujan Journal \textbf{6} (2002), 449--468.
Preprint arXiv:math/0107082 (2001).
\url{https://arxiv.org/abs/math/0107082}.

\bibitem{incoming}
ProveIt research continuation prepared with OpenAI assistance,
\emph{Gauss--Hurwitz Pole Cancellation: Resonant gamma-ratio sums,
harmonic identities, Stieltjes jets, and polylogarithmic endpoint
subtraction}, 10 October 2026, 23-page PDF. In particular, Section~9,
question~1, printed p.~19. Matching archive:
\file{docs/incoming/ProveIt_Gauss_Hurwitz_2026-10-10.zip},
Git blob \texttt{069b39cdc5a241c3e016f1780fd481dd3069d63a}.
Repository: \url{https://github.com/VladimirReshetnikov/ProveIt}.

\bibitem{repo-readme}
ProveIt contributors, \emph{Canonical polylogarithms manuscript README},
canonical manuscript README. Inspected 10 October 2026;
README blob \texttt{02774c39a68fff3792105aed3e2df11e3e604484}.
\url{https://github.com/VladimirReshetnikov/ProveIt/tree/main/Analysis/Polylogarithms/docs/manuscript}.

\bibitem{repo-integration}
ProveIt contributors, \emph{Zeta jets, Stieltjes antiderivatives and
negative polygamma}, canonical source
\file{Analysis/Polylogarithms/docs/manuscript/chapters/07-integration.tex}.
Inspected 10 October 2026. Repository:
\url{https://github.com/VladimirReshetnikov/ProveIt}.

\end{thebibliography}
\end{document}
